\documentstyle[% 


righttag% 


,11pt% 


,amssymb% 


,russian%               remove in Eng. 


,izvru%                 Set izven% in Eng. 


]{amsart} 


 


%


 


\newcommand{\eqdef}{\overset{\mathrm{def}}{=}{}} 


\newcommand{\const}{\operatorname{const}} 


\newcommand{\tts}[1]{} 


 


%\hoffset=-15mm%                  For Eng. 


\setcounter{page}{3} 


\god{2004} 


\nomer{5~(504)} %                        Remove in Eng. 


%\nomer{Vol.\,48, No.\,5}%               Set in Eng. 


%\str{\pageref{begin}--\pageref{end}}%  Set in Eng. 


\udk{514.756} 


 


\author{\it Д.А.\,Абруков} 


% NB:   ^^ remove in Eng. 


\title{Полярные гиперповерхности \\ 


проективно-метрического пространства} 


 


\date{} 


%\pagestyle{myheadings}%         Set in Eng. 


\pagestyle{plain} %              Remove in Eng. 


\begin{document} 


%\thispagestyle{plain}%          Set in Eng. 


\maketitle 


\label{begin} 


 


 


В работе изучаются некоторые вопросы геометрии двух регулярных 


гиперповерхностей $V_{n-1}$ и $\wt V_{n-1}$, погруженных в 


проективно-метрическое пространство $K_n$ и являющихся полярными 


относительно его абсолюта $Q_{n-1}$. 


 


Следующие индексы принимают значения 


\begin{gather*} 


\ov I,\ov K,\ov L=\ov{0,n};\quad I,K,L,P,Q=\ov{1,n};\\ 


i,j,k,p,q,s,l,t,u,v=\ov{1,n-1};\quad 


\ov i,\ov j,\ov k=\ov{0,n-1}. 


\end{gather*} 


 


Операция внешнего дифференцирования обозначена буквой $D$, 


внешнего умножения --- символом ``$\wedge$''. Оператор $\nabla$ действует по 


следующему закону [1]: 


$$ 


\nabla K^\alpha_{in}=d K^\alpha_{in}-K^\alpha_{tn} 


\omega^t_i-K^\alpha_{in}\omega^n_n+K^\beta_{in} 


\omega^\alpha_\beta; 


$$ 


при фиксированных главных параметрах этот оператор обозначается через 


$\nabla_\delta$, формы $\omega^{\ov I}_{\ov K}$ --- через $\pi^{\ov I}_{\ov 


K}$. Оператор $\wt\nabla$ действует по закону 


$$ 


\wt\nabla T^\alpha_{in}=dT^\alpha_{in}-T^\alpha_{tn} 


\Omega^t_i-T^\alpha_{in}\Omega^n_n+T^\beta_{in} 


\Omega^\alpha_\beta. 


$$ 


 


{\bf 1.} {\it Проективно-метрическое пространство}. Рассмотрим $n$-мерное 


проективное пространство $P_n$, отнесенное к подвижному реперу $R=\{A_{\ov 


K}\}$; 


деривационные формулы репера $R$ и уравнение структуры проективного 


пространства имеют соответственно вид [2]


\begin{alignat}{2} 


&\ \text{а)} &\quad & dA_{\ov I}=\omega^{\ov L}_{\ov I} 


A_{\ov L},\tag{1-a}\\ 


% 


&\ \text{б)} &\quad & D\omega^{\ov I}_{\ov K}= 


\omega^{\ov L}_{\ov K}\wedge\omega^{\ov I}_{\ov L},\quad 


\omega^{\ov L}_{\ov L}=0.\tag{1-б} 


\end{alignat} 


\addtocounter{equation}{1}


\noindent{}Известно [3], что проективно-метрическим пространством $K_n$ 


называется пространство $P_n$, в котором задана неподвижная гиперквадрика 


$Q_{n-1}$ (абсолют): 


\begin{equation} 


g_{\ov I\,\ov K}x^{\ov I}x^{\ov K}=0,\quad 


g_{\ov I\,\ov K}=g_{\ov K\,\ov I}, 


\end{equation} 


где $x^{\ov I}$ --- координаты точек $M\in Q_{n-1}$; следовательно, 


фундаментальной группой пространства $K_n$ является подгруппа группы 


проективных преобразований пространства $P_n$, а именно, стационарная 


подгруппа абсолюта $Q_{n-1}$. Согласно [4], условием неподвижности 


гиперквадрики (2) является выполнение дифференциальных уравнений 


\begin{equation} 


dg_{\ov I\,\ov K}-g_{\ov I\,\ov L}\omega^{\ov L}_{\ov K}- 


g_{\ov L\,\ov K}\omega^{\ov L}_{\ov I}=\Omega g_{\ov I\,\ov K}, 


\end{equation} 


где $\Omega$ --- некоторая форма Пфаффа. 


 


Известно [5], что за счет нормировки коэффициентов $g_{\ov I\,\ov K}$ 


гиперквадрики и вершин репера $R$ можно уравнение (2) абсолюта $Q_{n-1}$ и 


условие его неподвижности (3) записать соответственно в виде 


\begin{gather} 


a_{IK}x^I x^K+\frac{1}{c}(g_{I0}x^I+cx^0)^2=0,\\ 


% 


da_{IK}-a_{IL}\omega^L_K-a_{LK}\omega^L_I=-\frac{1}{c} 


(a_{IL}g_{K0}+a_{KL}g_{I0})\omega^L_0,\tag{5-а}\\ 


% 


dg_{I0}-g_{L0}\omega^L_I-c\omega^0_I=a_{IL}\omega^L_0,\tag{5-б} 


\end{gather} 


\addtocounter{equation}{1}





\noindent{}где $a_{IK}=g_{IK}-\frac{g_{I0}g_{K0}}{c}$, $a_{IK}=a_{KI}$, 


$c=g_{00}=\const\ne0$. 


Отметим, что при нормировке вершин репера $R$ получаем соотношение 


$\omega^0_0=-\frac{g_{0L}}{c}\omega^L_0$. 


В силу выражения (7) деривационные уравнения (1-a) в 


проективно-метрическом пространстве $K_n$ запишутся в виде 


$$ 


dA_0=\bigg(-\frac{g_{0L}}{c}A_0+A_L \bigg)\omega^L_0,\quad 


dA_I=\omega^{\ov L}_I A_{\ov L}. 


$$ 


Заметим также, что в силу соотношений (1), (7) имеет место 


$$ 


\omega^L_L=\frac{g_{0L}}{c}\omega^L_0. 


$$ 


Уравнения структуры (1-б) проективно-метрического пространства $K_n$ 


примут вид 


\begin{equation} 


\begin{split} 


& D\omega^I_0=-\frac{g_{0L}}{c}\omega^L_0\wedge \omega^I_0+ 


\omega^L_0\wedge\omega^I_L,\quad D\omega^0_I= 


\omega^L_I\wedge\omega^0_L-\frac{g_{0L}}{c}\omega^0_I 


\wedge\omega^L_0,\\ 


% 


&D\omega^0_0=\omega^L_0\wedge\omega^0_L,\quad 


D\omega^I_K=\omega^0_K\wedge\omega^I_0+ 


\omega^L_K\wedge\omega^I_L,\quad \omega^L_L= 


\frac{g_{0L}}{c}\omega^L_0. 


\end{split} 


\end{equation} 


 


{\bf 2.} {\it Дифференциальные уравнения гиперповерхности}. Как и в 


проективном пространстве $P_n$ [4], уравнениe гиперповерхности $V_{n-1}$ в 


$K_n$ в репере первого порядка ($A_0\in V_{n-1}$, $A_i\in T_{n-1}(A_0)$) 


записываeтся в виде 


\begin{equation} 


\omega^n_0=0. 


\end{equation} 


 


Отметим, что в силу $g_{00}\ne0$ текущая точка $A_0$ гиперповерхности 


$V_{n-1}$ не лежит на абсолюте. Трехкратное продолжение уравнения (7) 


приводит к следующим дифференциальным уравнениям компонент полей 


фундаментальных объектов второго $\{\Lambda^n_{ij}\}$ и третьего 


$\{\Lambda^n_{ij},\Lambda^n_{ijk}\}$ порядков $V_{n-1}\subset K_n$: 


\begin{gather} 


\omega^n_i=\Lambda^n_{ij}\omega^j_0,\quad \Lambda^n_{[ij]}=0,\\ 


% 


\nabla\Lambda^n_{ij}=\Lambda^n_{ijk}\omega^k_0,\quad 


\Lambda^n_{i[jk]}=\frac{1}{c}\Lambda{}^n_{i[j}g{}_{k]0},\quad 


\Lambda^n_{[ij]k}=0,\\ 


% 


\nabla\Lambda^n_{ijk}+\Lambda^n_{ik}\omega^0_j+ 


\Lambda^n_{jk}\omega^0_i-\Lambda{}^n_{(ij}\Lambda{}^n_{k)s} 


\omega^s_n=\Lambda^n_{ijkl}\omega^l_0,\quad 


\Lambda^n_{ij[kl]}=\frac{1}{c}\Lambda{}^n_{ij[k} 


g{}_{l]0}.\notag 


\end{gather} 


Из уравнений (9) видно, что система функций $\Lambda^n_{ij}$ образует 


симметричный тензор второго порядка. 


 


Заметим, что при $I=i$ из уравнений (5-б), (8) имеем 


$$ 


d\bigg(\frac{g_{i0}}{c} \bigg)-\frac{g_{j0}}{c}\omega^j_i- 


\omega^0_i=\wt a_{ij}\omega^j_0,\quad \wt a_{ij}=\frac{1}{c} 


(a_{ij}+g_{n0}\Lambda^n_{ij}). 


$$ 


 


Предполагая гиперповерхность $V_{n-1}$ в $K_n$ регулярной (т.\,е. 


$\Lambda=|\Lambda^n_{ij}|\ne0$), можно ввести обращенный тензор 


$\Lambda^{ij}_n$ второго порядка 


\begin{equation} 


\Lambda^{ik}_n\Lambda^n_{kj}=\delta^i_j,\quad 


\nabla\Lambda^{ij}_n=-\Lambda^{is}_n\Lambda^{tj}_n 


\Lambda^n_{stl}\omega^l_0. 


\end{equation} 


Функция $\Lambda$ есть относительный инвариант второго порядка 


$$ 


d\ln\Lambda+(n+1)\omega^n_n=A_l\omega^l_0,\quad 


A_l=\Lambda^{ji}_n\Lambda^n_{ijl}+\frac{2}{c}g_{l0}. 


$$ 


Продолжая последнее уравнение, имеем 


$$ 


\nabla A_i-(n+1)\Lambda^n_{ij}\omega^j_n=A_{ij}\omega^j_0,\quad 


A_{[ij]}=\frac{1}{c}A{}_{[i}g{}_{j]0}. 


$$ 


 


В четвертой дифференциальной окрестности точки $A_0\in V_{n-1}$ в $K_n$ 


внутренним образом определяется поле инвариантных соприкасающихся 


гиперквадрик $Q^2_{n-1}$ [4]: 


\begin{equation} 


\Lambda^n_{ij}x^i x^j+\frac{2p_i}{n+1}x^i x^n+S_n(x^n)^2= 


2x^0 x^n; 


\end{equation} 


здесь 


\begin{equation} 


p_i=\Lambda^{jk}_n\Lambda^n_{ijk}-\frac{g_{i0}}{c},\quad 


S_n=\frac{1}{n^2-1}\Lambda^{ij}_n\bigg(p_{ij}- 


\frac{p_i p_j}{n+1}-p_i\frac{g_{j0}}{c}\bigg). 


\end{equation} 


Функции $p_{ij}$ входят в дифференциальные уравнения 


\begin{gather*} 


\nabla p_i+(n+1)[\omega^0_i-\Lambda^n_{is}\omega^s_n]= 


p_{is}\omega^s_0,\\ 


\nabla_\delta p_{ij}+\bigg[p_j+(n+1)\frac{g_{j0}}{c}\bigg] 


\pi^0_i-[(n+1)a^n_{sij}+p_s\Lambda^n_{ij}]\pi^s_n+ 


2(n+1)\Lambda^n_{ij}\pi^0_n=0. 


\end{gather*} 


 


Аналогично [4] компоненты симметричного по всем нижним индексам 


тензора Дарбу на регулярной гиперповерхности $V_{n-1}$ в $K_n$ имеют 


вид 


\begin{align} 


b^n_{ijk}\eqdef&(n+1)\bigg(\Lambda^n_{ijk}+\frac{1}{c} 


\Lambda{}^n_{k(i}g{}_{j)0}\bigg)-\Lambda{}^n_{(ij}A{}_{k)},\quad 


\nabla b^n_{ijk}=b^n_{ijks}\omega^s_0,\\ 


% 


b^n_{ijks}=&(n+1)\bigg\{\Lambda^n_{ijks}+\frac{1}{c}[ 


g{}_{0(i}\Lambda{}^n_{j)ks}+\{a{}_{s(i}+g_{n0}\Lambda{}_{s(i}^n\} 


\Lambda{}^n_{j)k}]\bigg\}-\Lambda{}^n_{(ij|s|}A{}_{k)}- 


\Lambda{}^n_{(ij}A{}_{k)s}.\notag 


\end{align} 


 


Обращение в нуль тензора Дарбу регулярной гиперповерхности $V_{n-1}$ в 


$K_n$ есть условие ее вырождения в гиперквадрику $Q^2_{n-1}$ (см. (11)). 


\medskip 


 


{\bf 3.} {\it Полярная гиперповерхность}. Гиперповерхности $V_{n-1}$ и $\wt 


V_{n-1}$ назовем {\it полярными} (относительно абсолюта $Q_{n-1}$), если 


касательной гиперплоскостью $\wt T_{n-1}$ ($T_{n-1}$) в текущей точке 


$B\in\wt V_{n-1}$ ($A_0\in V_{n-1}$) будет поляра точки $A_0$ ($B$) 


относительно абсолюта $Q_{n-1}$ проективно-метрического 


пространства $K_n$. Нашей дальнейшей задачей является построение 


гиперповерхности $\wt V_{n-1}$, полярной данной гиперповерхности $V_{n-1}$ 


относительно абсолюта $Q_{n-1}$ пространства $K_n$. 


 


Пусть гиперповерхность $V_{n-1}$ в $K_n$ оснащена в смысле Э.\,Картана [6], 


т.\,е. каждой точке $A_0\in V_{n-1}$ присоединена точка 


\begin{equation} 


B_n(A_0)=v_n A_0+v^j_n A_j+A_n, 


\end{equation} 


не принадлежащая касательной плоскости $T_{n-1}(A_0)$; оснащающий объект 


$\{v^i_n,v_n\}$ подчинен условию инвариантности поля оснащающих точек: 


\begin{align*} 


dv^i_n-v^i_n\omega^n_n+v^j_n\omega^i_j+\omega^i_n&= 


v^i_{nk}\omega^k_0,\\ 


% 


dv_n-v_n\omega^n_n+v^j_n\omega^0_j+\omega^0_n&= 


v_{nk}\omega^k_0. 


\end{align*} 


 


Поле квазитензора $v^i_n$ определяет поле инвариантных нормалей первого 


рода гиперповерхности $V_{n-1}$ в $K_n$. 


 


Потребуем, чтобы точка $B_n$ принадлежала поляре $\pi$ точки $A_0$ 


относительно 


абсолюта $Q_{n-1}$: 


\begin{equation} 


\pi:g_{I0}x^I+c x^0=0. 


\end{equation} 


Так как условием принадлежности точки $B_n$ поляре (15) является выполнение 


равенства $g_{i0}v^i_n+c v_n+g_{n0}=0$, то справедливо 


\begin{equation} 


v_n=-\frac{1}{c}(g_{i0}v^i_n+g_{n0}); 


\end{equation} 


охват (16) определяет оснащающую точку Картана на нормали первого 


рода $v^i_n$. 


 


Полярой точки $B_n$ (см. (14)) относительно абсолюта $Q_{n-1}$ пространства 


$K_n$ является гиперплоскость 


\begin{equation} 


a_{Kj}x^K v^j_n+a_{Kn}x^K+\frac{1}{c}(g_{K0}x^K+c x^0) 


(g_{j0}v^j_n+g_{n0}+c v_n)=0. 


\end{equation} 


Соотношение (17) в силу равенств (16) запишется в виде 


\begin{equation} 


(a_{ij}v^j_n+a_{in})x^i+(a_{nj}v^j_n+a_{nn})x^n=0. 


\end{equation} 


Потребуем совпадения гиперплоскости (18) с касательной гиперплоскостью 


$T_{n-1}(A_0)$ гиперповерхности $V_{n-1}$ в $K_n$, т.\,е. с гиперплоскостью 


$x^n=0$. 


В силу последнего равенства и соотношений (18) имеем 


\begin{equation} 


a_{ij}v^j_n+a_{in}=0,\quad a_{nj}v^j_n+a_{nn}\ne0. 


\end{equation} 


Заметим, что 


\begin{equation} 


B_n\notin Q_{n-1}\ \Leftrightarrow \ a_{nn}+a_{in}v^i_n\ne0. 


\end{equation} 


Из уравнений (5) неподвижности абсолюта $Q_{n-1}$ проективно-метрического 


пространства $K_n$ с учетом уравнения (7) имеем 


\begin{gather} 


da_{ij}-a_{ik}\omega^k_j-a_{kj}\omega^k_i=\bigg\{ 


\Lambda{}^n_{s(i}a{}_{j)n}-\frac{1}{c}a{}_{s(i}g{}_{j)0}\bigg\}\omega^s_0,\\ 


%


da_{in}-a_{in}\omega^n_n-a_{kn}\omega^k_i-a_{ik}\omega^k_n= 


\bigg\{a_{nn}\Lambda^n_{is}-\frac{1}{c}(a_{is}g_{n0}+ 


a_{ns}g_{i0})\bigg\}\omega^s_0,\notag\\ 


%


da_{nn}-2a_{nn}\omega^n_n-2a_{kn}\omega^k_n=-\frac{2}{c} 


g_{n0}a_{ns}\omega^s_0,\\ 


%


dg_{i0}-g_{j0}\omega^j_i-c\omega^0_i=(a_{is}+g_{n0} 


\Lambda^n_{is})\omega^s_0,\notag\\ 


%


dg_{n0}-g_{j0}\omega^j_n-g_{n0}\omega^n_n-c\omega^0_n= 


a_{ns}\sigma^s_0. 


\end{gather} 


 


Предполагая тензор $a_{ji}$ невырожденным (т.\,е. $a\eqdef|a_{ij}|\ne0$), 


имеем 


поле взаимного тензора $a^{ij}$: 


\begin{equation} 


a^{ik}a_{kj}=\delta^i_j,\quad \nabla a^{ij}=\frac{1}{c}g_{s0} 


a{}^{s(i}\omega{}^{j)}_0-a^{is}a^{kj}a{}_{n(s}\Lambda{}^n_{k)l} 


\omega^l_0. 


\end{equation} 


Из соотношений (19) в силу (24) находим 


\begin{equation} 


v^i_n=-a^{ik}a_{kn}. 


\end{equation} 


 


Таким образом, из соотношений (16) и (25) следует, что точка $B_n\in\pi$ 


(см. (14)), являющаяся полюсом касательной гиперплоскости $T_{n-1}(A_0)$ 


гиперповерхности $V_{n-1}$ в $K_n$ относительно абсолюта $Q_{n-1}$, в репере 


первого порядка $R=\{A_{\ov K}\}$ имеет разложение 


\begin{equation} 


B_n=\frac{1}{c}(g_{j0}a_{sn}a^{sj}-g_{n0})A_0-a_{sn}a^{sj} 


A_j+A_n. 


\end{equation} 


 


Отметим, что условие (20) с учетом соотношений (25) перепишется в виде 


\begin{equation} 


B_n\notin Q^2_{n-1}\ \Leftrightarrow\ A_{nn}\eqdef a_{nn}-a_{in}a_{jn}a^{ij}\ne0, 


\end{equation} 


где $A_{nn}$ --- относительный инвариант, удовлетворяющий уравнению 


\begin{equation} 


dA_{nn}-2A_{nn}\omega^n_n=-2A_{nn}a^{ij}a_{in}\Lambda^n_{jk}\omega^k_0. 


\end{equation} 


 


Итак, исходная гиперповерхность $V_{n-1}$ в $K_n$ внутренним образом 


индуцирует 


гиперповерхность $\wt V_{n-1}$, полярную данной $V_{n-1}$ относительно 


абсолюта $Q_{n-1}$. Касательной гиперплоскостью $\wt T_{n-1}$ в текущей точке 


$B_n\in\wt V_{n-1}$ (см. (26)) будет поляра $\pi$ точки $A_0$. 


Гиперповерхности $V_{n-1}$ и $\wt V_{n-1}$ являются полярными относительно 


абсолюта $Q_{n-1}$. 


 


Так как $A_0\notin\pi$, то, рассматривая точки $B_i=A_i+x^0_i A_0$, 


потребуем, чтобы они принадлежали поляре $\pi$ точки $A_0$ (см. (15)): 


$g_{i0}+c x^0_i=0$. Отсюда находим $x^0_i=-\frac{1}{c}g_{i0}$. Следовательно, 


\begin{equation} 


B_i=A_i-\frac{1}{c}g_{i0}A_0. 


\end{equation} 


 


Таким образом, с полярной гиперповерхностью $\wt V_{n-1}$ связан текущий 


репер первого порядка $\wt R=\{B_n,B_i,B_0\}=\{B_{\ov I}\}$, где $B_n$ и 


$B_i$ имеют соответственно строения (26) и (29), а в качестве вершины $B_0$ 


примем $A_0$: $B_0\equiv A_0$. 


 


Деривационные формулы репера $\wt R=\{B_{\ov K}\}$ и уравнения структуры 


проективно-метрического пространства $K_n$ (в репере $\wt R$) запишутся 


соответственно в виде 


\begin{equation} 


\begin{split} 


dB_{\ov I}&=\Omega^{\ov L}_{\ov I}B_{\ov L},\\ 


D\Omega^{\ov I}_{\ov K}&=\Omega^{\ov L}_{\ov K}\wedge\Omega^{\ov I}_{\ov L}, 


\end{split} 


\end{equation} 


где формы $\Omega^{\ov K}_{\ov I}$ имеют строения 


\begin{gather} 


\Omega^n_0=\omega^n_0=0,\quad \Omega^0_0=0,\quad 


\Omega^n_n=\omega^n_n-a_{sn}a^{sj}\omega^n_j,\notag\\ 


\Omega^n_i=\omega^n_i,\quad \Omega^0_i=-\frac{1}{c} 


a_{ij}\omega^j_0,\quad \Omega^i_n=-A_{nn}a^{is}\omega^n_s,\\ 


\Omega^i_j=\omega^i_j+a^{is}a_{sn}\omega^n_j-\frac{1}{c} 


g_{j0}\omega^i_0,\quad \Omega^0_n=-\frac{A_{nn}}{c}\omega^n_0=0.\notag 


\end{gather} 


 


Отметим, что в силу соотношений (6) и (31) имеем $\Omega^{\ov L}_{\ov L}=0$. 


В силу соотношений (30), (31) система форм $\{\Omega^k_n\}$ на 


гиперповерхности 


$\wt V_{n-1}$ будет базисной: 


$\Omega^1_n\wedge\dotsb\wedge\Omega^{n-1}_n\ne0$. 


Уравнение полярной гиперповерхности $\wt V_{n-1}$ в $K_n$ в репере первого 


порядка $\wt R=\{B_{\ov K}\}$ имеет вид 


\begin{equation} 


\Omega^0_n=0. 


\end{equation} 


Заметим, что согласно соотношениям (31), (32) справедливо 


\begin{equation} 


\Omega^n_i=-\frac{1}{A_{nn}}a_{is}\Omega^s_n,\quad 


\Omega^i_0=-\frac{1}{A_{nn}}\Lambda^{il}_n a_{ls}\Omega^s_n. 


\end{equation} 


 


Трехкратное продолжение уравнений (32) с использованием соотношений (33) 


приводит к следующим дифференциальным уравнениям компонент полей 


фундаментальных объектов второго $\{V^n_{ij}\}$ и третьего $\{V^n_{ijk}\}$ 


порядков гиперповерхности $\wt V_{n-1}$ в $K_n$: 


\begin{gather} 


\Omega^0_i=V^n_{ij}\Omega^j_n,\quad V^n_{[ij]}=0;\notag\\ 


%


\wt\nabla V^n_{ij}+V^n_{ij}\Omega^0_0=V^n_{ijk}\Omega^k_n, 


\quad V^n_{i[jk]}=0;\\ 


%


dV^n_{ijk}+2V^n_{ijk}\Omega^n_n-V{}^n_{s(ij}\Omega{}_{k)}^s= 


V^n_{ijks}\Omega^s_n,\notag\\ 


%


V^n_{ij[ks]}=\frac{1}{A_{nn}}\{V{}_{i[k}^n a{}_{s]j}+ 


V{}^n_{j[k}a{}_{s]i}-\Lambda^{lt}_n(V^n_{il}V{}^n_{j[k} 


a{}_{s]t}+V^n_{jl}V{}^n_{i[k}a{}_{s]t})\}\notag. 


\end{gather} 


Отметим, что согласно уравнениям (34) и соотношениям (31) имеем 


\begin{equation} 


V^n_{ij}=\frac{1}{c A_{nn}}\Lambda^{lt}_n a_{il}a_{jt}. 


\end{equation} 


 


Заметим, что с учетом последних соотношений выражения (33) перепишутся 


в виде $\Omega^n_i=-\frac{1}{A_{nn}}a_{is}\Omega^s_n$, $\Omega^i_0=-c 


a^{it}V^n_{ts}\Omega^s_n$. 


 


Дифференцируя выражения (35), с использованием уравнений (10), 


(21), (28), (34) и соотношений (31) получим 


\begin{equation} 


V^n_{ijk}=\frac{1}{c A^2_{nn}}\{\Lambda^{sp}_n 


\Lambda^{qt}_n\Lambda^{lu}_n\Lambda^n_{pql}a_{is}a_{jt}a_{ku}- 


c A_{nn}V{}^n_{(ij}a{}_{k)n}+A_{nn}\Lambda^{st}_n g_{t0} 


a{}_{s(i}v{}^n_{j)k}\}. 


\end{equation} 


 


Из уравнений (34) видно, что каждая из систем функций $\{V^n_{ij}\}$, 


$\{V^n_{ijk}\}$ образует симметричный тензор второго и третьего порядка 


соответственно. 


 


Согласно выражениям (35) справедлива 


 


\begin{thm} 


При сделанных выше предположениях $(10)$, $(24)$, $(27)$ полярная 


гиперповерхность $\wt V_{n-1}$ в $K_n$ является регулярной, т.\,е. 


$V\eqdef|v^n_{ij}|\ne0$. 


\end{thm} 


 


В силу регулярности полярной гиперповерхности $\wt V_{n-1}$ в $K_n$ можно 


ввести обращенный тензор $V^{ij}_n$ второго порядка $V^{ik}_n 


V^n_{kj}=\delta^i_j$, $\wt \nabla V^{ij}_n=-V^{is}_n V^{tj}_n V^n_{stl}\Omega^l_n$, 


причем 


\begin{equation} 


V^{ij}_n=c A_{nn} a^{ik}a^{jt}\Lambda^n_{kt}. 


\end{equation} 


 


Функция $V$ есть относительный инвариант второго порядка 


$$ 


d\ln V+(n+1)\Omega^n_n=V_k\Omega^k_n, 


$$ 


где 


\begin{equation} 


V_k=V^{ji}_n V^n_{ijk}. 


\end{equation} 


Продолжая последнее уравнение, имеем 


$$ 


\wt\nabla V_i+V_i\Omega^n_n=V_{ij}\Omega^j_n,\quad V_{[ij]}=0; 


$$ 


следовательно, совокупность функций $V_i$ образует тензор третьего порядка. 


 


Заметим, что согласно соотношениям (38) с использованием равенств (36) 


и (37) имеем 


\begin{equation} 


V_k=\frac{1}{A_{nn}}(\Lambda^{st}_n A_s a_{kt}-(n+1)a_{kn}). 


\end{equation} 


 


Уравнение гиперквадрики (абсолюта) $Q_{n-1}$ проективно-метрического 


пространства $K_n$ в репере $R=\{A_{\ov K}\}$ имеет вид (4). Запишем данное 


уравнение в репере $\wt R=\{B_{\ov K}\}$ (см. (26), (29)). 


 


Известно [7], что связь между старыми $x^{\ov I}$ и новыми $y^{\ov 


I}$ проективными координатами точки $M\in Q_{n-1}$ имеет вид 


$$ 


\rho x^{\ov I}=c^{\ov I}_{\ov K}y^{\ov K}, 


$$ 


где $\ov K$-й столбец матрицы $c^{\ov I}_{\ov K}$ состоит из координат вершин 


$B_{\ov K}$ нового репера $\wt R$ в старом репере $R$. В силу (26), (29) и 


$B_0\equiv A_0$ имеем 


$$ 


c^{\ov I}_{\ov K}= 


\begin{pmatrix} 


1 & -\frac{g_{j0}}{c}& \frac{1}{c}\{a^{kt}g_{k0}a_{tn}-g_{n0}\}\\ 


0 & \delta^i_j & -a^{ki}a_{kn}\\ 


0 & 0 & 1 


\end{pmatrix}. 


$$ 


 


Таким образом, связь между старыми $x^{\ov I}$ и новыми $y^{\ov I}$ 


координатами 


точки $M\in Q_{n-1}$ имеет вид 


\begin{align*} 


x^0&=y^0-\frac{g_{j0}}{c}y^j+\frac{1}{c}\{a^{kt}g_{k0} 


a_{tn}-g_{n0}\}y^n,\\ 


x^i&=y^i-a^{ik}a_{kn}y^n,\\ 


x^n&=y^n. 


\end{align*} 


Подставляя последние соотношения в уравнение (4), получаем уравнение 


абсолюта $Q_{n-1}$ в новом репере $\wt R=\{B_n,B_i,B_0\}=\{B_{\ov I}\}$: 


\begin{equation} 


a_{ij}y^i y^j+A_{nn}(y^n)^2+c(y^0)^2=0, 


\end{equation} 


где $y^{\ov K}$ --- координаты точек $\wt M\in Q_{n-1}$ в репере $\wt R$. 


 


Согласно уравнениям (21)--(22), (24), выражению $A_{nn}$ (см. (27)) и 


соотношениям (31) находим 


\begin{equation} 


da_{ij}-a_{ik}\Omega^k_j-a_{kj}\Omega^k_i=0,\quad 


dA_{nn}-2A_{nn}\Omega^n_n=0. 


\end{equation} 


Уравнения (41) являются (в репере $\wt R$) условием неподвижности абсолюта 


$Q_{n-1}$ (см. (40)). 


 


Вернемся к полярной гиперповерхности $\wt V_{n-1}$. Следуя [4], в 


четвертой дифференциальной окрестности точки $B_n\in\wt V_{n-1}$ внутренним 


образом определяется поле инвариантных соприкасающихся гиперквадрик $\wt 


Q^2_{n-1}$: 


\begin{equation} 


V^n_{ij}y^i y^j+\frac{2V_i}{n+1}y^i y^0+S_0(y^0)^2=2y^0y^n, 


\end{equation} 


где 


$$ 


S_0=\frac{1}{n^2-1}V^{ij}_n\bigg(V_{ij}-\frac{V_i V_j}{n+1} \bigg)+ 


\frac{1}{n-1}\bigg(c V^n_{ij}a^{ij}-\frac{1}{A_{nn}}V^{ij}_n a_{ij} 


\bigg). 


$$ 


 


Компоненты симметричного тензора Дарбу на регулярной гиперповерхности 


$\wt V_{n-1}$ в $K_n$ имеют вид 


\begin{equation} 


D^n_{ijk}\eqdef(n+1)V^n_{ijk}-V{}^n_{(ij}V{}_{k)},\quad 


\wt\nabla D^n_{ijk}+D^n_{ijk}\Omega^n_n=D^n_{ijks}\Omega^s_n. 


\end{equation} 


 


Аналогично работе [4] справедливо следующее: необходимым и достаточным 


условием вырождения регулярной гиперповерхности $\wt V_{n-1}$ в $K_n$ в 


гиперквадрику $\wt Q^2_{n-1}$ (см. (42)) является обращение в нуль 


симметричного тензора Дарбу (43). 


 


Согласно соотношениям (13), (35), (36), (39), (43) компоненты 


симметричных тензоров Дарбу полярных регулярных гиперповерхностей $V_{n-1}$ 


и $\wt V_{n-1}$ в $K_n$ связаны равенствами 


\begin{equation} 


D^n_{ijk}=\frac{1}{c A^2_{nn}}\Lambda^{sp}_n 


\Lambda^{qt}_n\Lambda^{lv}_n a_{is}a_{jt}a_{kv}b^n_{pql}. 


\end{equation} 


 


Таким образом, доказана 


 


\begin{thm} 


Гиперповерхность $V_{n-1}$ в $K_n$ вырождается в гиперквадрику $Q^2_{n-1}$ 


$($см. $(11))$ тогда и только тогда, когда полярная ей $($относительно 


абсолюта $Q_{n-1})$ гиперповерхность $\wt V_{n-1}$ вырождается в 


гиперквадрику $\wt Q^2_{n-1}$ $($см. $(42))$. 


\end{thm} 


 


{\bf 4.} {\it Полярные нормализации гиперповерхностей $V_{n-1}$ и $\wt 


V_{n-1}$ в $K_n$}. Оснащение в смысле А.П.\,Нордена [3] гиперповерхности 


$V_{n-1}$ в $K_n$ равносильно заданию на этом подмногообразии двух полей 


квазитензоров $v^i_n$, $v_i$: 


\begin{equation} 


\nabla v^i_n+\omega^i_n=v^i_{nj}\omega^j_0,\quad 


\nabla v_i+\omega^0_i=v_{ij}\omega^j_0, 


\end{equation} 


определяющих соответственно поля нормалей первого $N_1\equiv[A_0,N_n]$ и 


второго $N_{n-2}\equiv[N_i]$ родов на гиперповерхности $V_{n-1}$ в $K_n$; 


здесь $N_n=A_n+v^i_n A_i$, $N_i=A_i+v_i A_0$. 


 


Для системы функций 


\begin{equation} 


E_i(v)=\frac{p_i}{n+1}-v_i+\Lambda^n_{is}v^s_n 


\end{equation} 


в силу уравнений (9), (12), (45) имеем 


$$ 


\nabla E_i=E_{ij}\omega^j_0,\quad E_{[ij]}= 


\frac{A{}_{[i}g{}_{j]0}}{c(n+1)}-v_{[ij]}+ 


\frac{v^s_n\Lambda{}^n_{s[i}g{}_{j]0}}{c}-v{}^s_{n[i}\Lambda{}^n_{j]s}. 


$$ 


 


Обращение в нуль тензора $E_i(v)$ есть условие, при котором нормализация 


$\{v^i_n,v_i\}$ гиперповерхности $V_{n-1}$ в $K_n$ является взаимной [3] 


относительно поля соприкасающихся гиперквадрик (11). 


 


Поля квазитензоров четвертого порядка [1]


\begin{equation} 


F^i_n=\frac{1}{2}\Lambda^{ik}_n\bigg(C_k-\frac{p_k}{n+1} 


-\frac{g_{k0}}{c} \bigg),\quad F_i=\frac{1}{2}\bigg(C_i+ 


\frac{p_i}{n+1}-\frac{g_{i0}}{c}\bigg) 


\end{equation} 


на гиперповерхности $V_{n-1}$ в $K_n$ определяют инвариантные нормали Фубини 


первого и второго родов соответственно. 


 


Отметим, что функции $C_k$ входят в дифференциальное уравнение 


\begin{equation} 


d\ln C_n-\omega^n_n=C_k\omega^k_0 


\end{equation} 


ненулевого относительного инварианта $C_n$ третьего порядка 


$$ 


C_n=\Lambda^{ij}_n\Lambda^{st}_n\Lambda^{kl}_n 


b^n_{isk}b^n_{jtl}. 


$$ 


Дифференциальные уравнения функций $C_i$ имеют вид $\nabla 


C_i+\Lambda^n_{ij}\omega^j_n=C_{ij}\omega^j_0$. 


 


Нормализация Фубини $(F^i_n,F_i)$ регулярной гиперповерхности $V_{n-1}$ 


является взаимной относительно поля соприкасающихся гиперквадрик 


(11), т.\,к. $E_i(F)=0$. 


 


Оснащение в смысле А.П.\,Нордена [3] полярной гиперповерхности 


$\wt V_{n-1}$ в $K_n$ равносильно заданию на этом многообразии двух полей 


тензоров $v^i_0$, $v^n_i$: 


$$ 


dv^i_0-v^i_0\Omega^0_0+v^t_0\Omega^i_t=v^i_{0j}\Omega^j_n,\quad 


dv^n_i-v^n_t\Omega^t_i+v^n_i\Omega^n_n=v^n_{ij}\Omega^j_n, 


$$ 


определяющих соответственно поля нормалей первого $\wt N_1\equiv[B_n,\wt 


N_0]$ и второго $\wt N_{n-2}\equiv[\wt N_i]$ родов на гиперповерхности $\wt 


V_{n-1}$; здесь $\wt N_0=B_0+v^i_0 B_i$, $\wt N_i=B_i+v^n_i B_n$. 


 


Нормаль первого рода гиперповерхности $V_{n-1}$ инцидентна точкам $A_0$ 


и $N_n=A_n+v^i_n A_i$. Поляры этих точек относительно абсолюта $Q_{n-1}$ (см. 


(4)) определяются соответственно уравнениями 


$$ 


g_{I0}x^I+cx^0=0,\quad a_{Kj}x^K v^j_n+a_{Kn}x^K+\frac{1}{c} 


(g_{K0}x^K+cx^0)(g_{j0}v^j_n+g_{n0})=0, 


$$ 


т.\,е. поляра прямой $[A_0,N_n]$ есть $(n-2)$-мерная плоскость, принадлежащая 


касательной гиперплоскости полярной гиперповерхности $\wt V_{n-1}$ в 


соответствующей точке $B_n\in\wt V_n$; уравнение этой $(n-2)$-мерной 


плоскости 


определяется системой 


\begin{equation} 


g_{I0}x^I+cx^0=0,\quad a_{Ik}x^I v^k_n+a_{In}x^I=0. 


\end{equation} 


 


Так как нормаль второго рода гиперповерхности $\wt V_{n-1}$ определяется 


точками $\wt N_i=B_i+v^n_i B_n$, то согласно соотношениям (26) и (29) 


получаем 


$$ 


\wt N_i=\bigg(-\frac{1}{c}g_{i0}+\frac{1}{c}v^n_i\{g_{t0}a_{sn}a^{st}


-g_{n0}\}\bigg)A_0+(\delta^j_i-a^{js}a_{sn}v^n_i)A_j+v^n_i A_n. 


$$ 


Потребуем, чтобы $(n-2)$-мерная плоскость (49) задавала нормаль второго 


рода гиперповерхности $\wt V_{n-1}$; для этого координаты точек $\wt N_i$ 


должны удовлетворять системе уравнений (49), что имеет место в случае 


выполнения соотношений 


\begin{equation} 


v^n_i=-\frac{1}{A_{nn}}(a_{ik}v^k_n+a_{in}). 


\end{equation} 


 


Таким образом, при полярном (относительно абсолюта $Q_{n-1}$) отображении 


нормаль первого рода $v^i_n$ гиперповерхности $V_{n-1}$ в $K_n$ переходит в 


нормаль второго рода $v^n_i$ полярной гиперповерхности $\wt V_{n-1}$ в $K_n$, 


определяемую квазитензором (50). 


 


Аналогично, нормаль второго рода $v_i$ гиперповерхности $V_{n-1}$ при 


полярном отображении переходит в нормаль первого рода $v^i_0$ полярной 


гиперповерхности $\wt V_{n-1}$, причем 


\begin{equation} 


v^i_0=-a^{ik}(g_{k0}+c v_k). 


\end{equation} 


 


Следовательно, справедлива 


 


\begin{thm} 


Нормализация в смысле А.П.\,Нордена одной из полярных 


гиперповерхностей $V_{n-1}$ или $\wt V_{n-1}$ в $K_n$ равносильна 


нормализации 


другой$;$ при этом нормаль первого $($второго$)$ рода $v^i_n$ $(v_i)$ 


гиперповерхности $V_{n-1}$ полярна $($относительно абсолюта $Q_{n-1})$ нормали 


второго $($первого$)$ рода $v^n_i$ $(v^i_0)$ гиперповерхности $\wt V_{n-1}$, 


причем оснащающие объекты $(v^i_n,v_i)$ и $(v^i_0,v^n_i)$ связаны 


соотношениями $(50)$, $(51)$. 


\end{thm} 


 


Нормализации взаимных гиперповерхностей $V_{n-1}$ и $\wt V_{n-1}$ в $K_n$ 


полями объектов $(v^i_n,v_i)$ и $(v^i_0,v^n_i)$ соответственно, связанные 


между собой соотношениями (50) и (51), назовем {\it полярными} по отношению 


друг к другу. 


 


Рассмотрим систему функций 


\begin{equation} 


\wt E_i(v)=\frac{V_i}{n+1}-v^n_i+V^n_{is}v^s_0,\quad 


\wt\nabla\wt E_i=\wt E_{ij}\Omega^j_n,\quad \wt E_{[ij]}= 


-v^n_{[ij]}-v{}^s_{0[j}V{}^n_{j]s}. 


\end{equation} 


Обращение в нуль тензора $\wt E_(v)$ есть условие, при котором нормализация 


$(v^i_0,v^n_i)$ гиперповерхности $\wt V_{n-1}$ в $K_n$ является взаимной 


относительно поля соприкасающихся гиперквадрик (42). 


 


Заметим, что согласно соотношениям (35), (39), (46), (50)--(52) 


получим $\wt E_i(v)=\frac{1}{A_{nn}}\Lambda^{st}_n a_{it}E_s(v)$. 


 


Таким образом, имеет место 


 


\begin{thm} 


Нормализация регулярной гиперповерхности $V_{n-1}$ в $K_n$ 


взаимна относительно поля соприкасающихся гиперквадрик $Q^2_{n-1}$ $($см. 


$(11))$ тогда и только тогда, когда взаимна относительно поля 


соприкасающихся гиперквадрик $\wt Q^2_{n-1}$ $($см. $(42))$ полярная 


относительно абсолюта $Q_{n-1}$ нормализация регулярной гиперповерхности $\wt 


V_{n-1}$. 


\end{thm} 


 


В четвертой дифференциальной окрестности определяются поля нормалей 


Фубини первого и второго родов 


\begin{equation} 


\Phi^i_0\eqdef\frac{1}{2}V^{ik}_n\bigg(\wt C_k-\frac{V_k}{n+1}\bigg), 


\quad \Phi^n_i\eqdef\frac{1}{2}\bigg(\wt C_i+\frac{V_i}{n+1}\bigg), 


\end{equation} 


где функции $\wt C_i$ входят в дифференциальное уравнение 


\begin{equation} 


d\ln\wt C_n+\Omega^n_n=\wt C_i\Omega^i_n 


\end{equation} 


ненулевого относительного инварианта 


\begin{equation} 


\wt C_n=V^{ij}_n V^{st}_n V^{kl}_n D^n_{isk}D^n_{jtl}. 


\end{equation} 


Дифференциальные уравнения функций $\wt C_i$ имеют вид $d\wt C_i-\wt C_j 


\Omega^j_i+\wt C_i\Omega^n_n=\wt C_{ij}\Omega^j_n$. 


 


Согласно соотношениям (50) и (51) можно получить поля нормалей 


$F^i_0$ и $F^n_i$ полярной гиперповерхности $\wt V_{n-1}$ в $K_n$, полярные 


относительно абсолюта (4) полям нормалей Фубини гиперповерхности $V_{n-1}$ в 


$K_n$ (см. (47)): 


\begin{equation} 


F^i_0=-a^{ik}(g_{k0}+c F_k),\quad F^n_i=-\frac{1}{A_{nn}} 


(a_{ik}F^k_n+a_{in}). 


\end{equation} 


 


Согласно равенствам (31), (35), (44), (48), (54), (55) имеем 


\begin{equation} 


\wt C_i=-\frac{1}{A_{nn}}(C_j\Lambda^{jk}_n a_{ik}+a_{in}). 


\end{equation} 


Рассматривая соотношения (35), (39), (53), (56), (57), получим 


$F^i_0=\Phi^i_0$, $F^n_i=\Phi^n_i$. 


 


Следовательно, доказана 


 


\begin{thm} 


Нормализация гиперповерхности $\wt V_{n-1}$ в $K_n$, полярная 


$($относительно абсолюта $Q_{n-1}$ пространства $K_n)$ нормализации Фубини 


исходной гиперповерхности $V_{n-1}$, является ее нормализацией Фубини 


полями нормалей $(53)$ или, что то же самое, полями нормалей $(56)$. 


\end{thm} 


 


Согласно теореме 4 нормализация Фубини $(F^i_0,F^n_i)$ регулярной 


гиперповерхности $\wt V_{n-1}$ в $K_n$ является взаимной относительно поля 


соприкасающихся гиперквадрик (52). 


 


\begin{thebibliography}{9} 


 


\bibitem{1} 


Столяров А.В. {\em Двойственная теория оснащенных многообразий}. -- 


Чебоксары, 1994. -- 290~с. 


\bibitem{2} 


Фиников С.П. {\em Метод внешних форм Картана в дифференциальной 


геометрии}. -- М.--Л.: ГИТТЛ, 1948. -- 432~с. 


\bibitem{3} 


Норден А.П. {\em Пространства аффинной связности}. -- М.: Наука, 


1976. -- 432~с. 


\bibitem{4} 


Лаптев Г.Ф. {\em Дифференциальная геометрия погруженных многообразий. 


Теоретико-групповой метод дифференциально-геометрических 


исследований} // Тр. Московск. матем. о-ва. -- 1953. 


-- Т.\,2. -- С.\,275--382. 


\bibitem{5} 


Столяров А.В. {\em Внутренняя геометрия проективно-метрического 


пространства} // Дифференц. геометрия многообразий фигур. -- 


Кaлининград, 2001. -- Вып.\,32. -- С.\,94--101. 


\bibitem{6} 


Cartan E. {\em Les espaces \'a connexion protective} // Тр. семин. по 


векторн. и тензорн. анализу. -- МГУ. -- 1937. -- Вып.\,4. -- С.\,147--159. 


 


\bibitem{7} 


Ефимов Н.В. {\em Высшая геометрия}. -- М.: ГИТТЛ, 1961. -- 580~с. 


 


\end{thebibliography} 


 


\vskip 2cm 


 


\postup{\it Чувашский государственный}{\qquad \it Поступили} 


\postup{\it педагогический университет}{\it первый вариант $28.11.2001$} 


\postup{}{\it окончательный вариант $13.10.2002$} 


 


\label{end} 


\end{document} 


